\chapter{Introduction}
\label{ch:introduction}

\section{Background}

Human activity recognition underpins applications from elderly-care
monitoring to smart-building automation. The established sensing modalities
each carry a cost that limits deployment. Cameras provide rich information
but raise serious privacy objections in exactly the private spaces --- homes,
bedrooms, bathrooms --- where fall detection matters most, and they degrade
in darkness and under occlusion. Wearable inertial sensors avoid the privacy
problem but depend on compliance: a device that is not worn, or whose battery
has run flat, detects nothing, and the elderly users who benefit most are the
least reliable wearers.

WiFi sensing offers a third option. A radio signal propagating from a
transmitter to a receiver arrives along many paths, reflecting and
diffracting around walls, furniture, and people. A moving body changes the
lengths and attenuations of those paths, so the received signal carries a
record of motion in the environment. Because ordinary WiFi hardware is
already deployed in most indoor spaces, this sensing is contactless,
device-free, works in complete darkness and through soft partitions, and
captures no imagery.

Two quantities are available. The Received Signal Strength Indicator (RSSI)
is a single scalar per packet describing total received power; it is coarse,
and its variation confounds motion with many other effects. Channel State
Information (CSI) is far richer: an OFDM receiver estimates the complex
channel response of every subcarrier separately, so a single 20\,MHz packet
yields tens of independent frequency-resolved measurements. Different
subcarriers are affected differently by the same motion, and it is this
frequency diversity that makes fine-grained activity classification
feasible.

Historically CSI extraction required specialised network cards --- the Intel
5300 or Atheros chipsets --- whose cost and driver requirements confined the
field to laboratories. Modern low-cost microcontrollers with integrated WiFi,
in particular the Espressif ESP32 family, now expose CSI directly, making
the sensing hardware for a multi-receiver system inexpensive enough for
deployment in cost-constrained settings.

\section{Problem Statement}

This work addresses the recognition of five activities --- walking, sitting,
standing, lying, and falling --- plus a background class denoting an empty
room, from WiFi CSI captured on commodity microcontroller hardware.

Falling is the safety-critical class. A missed fall is the expensive error,
and its cost is not symmetric with a false alarm; it is also the rarest class,
occurring as brief discrete events rather than sustained minutes of
behaviour. The background class is equally necessary for a credible system:
a classifier trained only on activities will confidently report an activity
in an empty room.

Two constraints shape the technical approach. First, the receivers have a
single antenna and a consumer-grade oscillator, so raw CSI phase is corrupted
by carrier and sampling frequency offsets, and the multi-antenna phase
sanitisation used with specialised cards is structurally unavailable. The
pipeline is therefore amplitude-only. Second, and more consequential for how
results should be read, activity recognition from CSI is only useful if it
generalises to a new day and a new person. A model that has memorised the
radio signature of one recording session is worthless in deployment while
appearing excellent under the evaluation protocol most commonly reported in
this literature.

\section{Objectives}
\begin{itemize}
  \item Build a 3-receiver ESP32-S3 CSI acquisition system (100 packets/s).
  \item Collect a labeled multi-subject, multi-session activity dataset.
  \item Train and compare classical and deep models (SVM/RF baseline,
        CNN, CNN-LSTM) on amplitude spectrograms.
  \item Achieve $\geq 85\%$ accuracy on random and cross-session splits;
        quantify the cross-subject generalization gap honestly.
  \item Demonstrate real-time recognition with temporal smoothing and a
        live dashboard, including a fall alert.
\end{itemize}

\section{Scope and Limitations}

The system is scoped to a single person in the sensing area at a time.
Multiple simultaneous occupants superimpose their perturbations on the same
channel, and separating them is a substantially harder problem that is not
attempted here.

Recognition covers the five activities and background listed above.
Transitional movements such as sitting down and standing up are not separate
classes; they are trimmed from the labelled spans rather than assigned to the
activity they precede or follow.

Sensing is confined to indoor rooms with the transmitter and receivers in
fixed, measured positions. Through-wall operation at range is not a target,
though it has been demonstrated on comparable hardware elsewhere.

Only CSI amplitude is used, for the hardware reasons given above. Phase
features are not built into the pipeline and are not claimed.

Finally, the models are trained and evaluated on data collected by this
project. Cross-environment generalisation is measured to the limited extent
that a second room permits, but a system deployed in an arbitrary unseen
room should be expected to require adaptation data.

\section{Report Organization}

Chapter~\ref{ch:literature} reviews WiFi sensing fundamentals, deep learning
approaches to CSI-based recognition, and prior work on ESP32-class hardware,
identifying the gap this project addresses.
Chapter~\ref{ch:methodology} sets out the acquisition, preprocessing,
modelling, and --- centrally --- the three-split evaluation protocol.
Chapter~\ref{ch:system} describes the hardware deployment, software
architecture, real-time inference path, and quality assurance.
Chapter~\ref{ch:results} presents the dataset, model comparison across
splits, fall detection performance, ablation studies, real-time measurements,
and failure analysis.
Chapter~\ref{ch:conclusion} concludes and identifies future work.
